% 《大学物理简明教程》第5版第10章解答题10.3.17—10.3.23；
% 原书与PDF第274页。光栅级数和缺级严格区分；图10.3.23按Fresnel规律提出核验基准。
\subsection{解答题：光栅缺级、偏振与布儒斯特定律}

\sourceproblem{10.3.17}{274}{274}
\problemtext 波长$\lambda=600$ nm的单色光垂直入射光栅，
第2、3级主明纹分别出现在
$\sin\varphi_2=0.20$、$\sin\varphi_3=0.30$，
且第4级主明纹缺级。求：(1)光栅常量；
(2)透光缝宽；(3)在$0^\circ<\varphi<90^\circ$内能观察到哪些主明纹级数。
\workings
光栅主极大条件
$d\sin\varphi_k=k\lambda$给
\[
\boxed{d=\frac{2(600\,\mathrm{nm})}{0.20}
=\frac{3(600\,\mathrm{nm})}{0.30}=6.00\,\mu\mathrm m}.
\]
第4级恰与单缝第1个零点相重合，
即$d/a=4$，因此
\[
\boxed{a=\frac d4=1.50\,\mu\mathrm m}.
\]
可传播的正级数满足$k<d/\lambda=10$
（严格排除$90^\circ$掠射方向）。
因此$k=1,\ldots,9$，又因$d/a=4$的倍数缺级，
$k=4,8$消失。剩余
\[
\boxed{k=1,2,3,5,6,7,9}.
\]
\crosscheck 题设给的是第2、3级\emph{主明纹}，
不能误按单缝第2、3级暗纹求$d$；
第4级缺级是单缝包络零点与光栅主极大重合。

\sourceproblem{10.3.18}{274}{274}
\problemtext 双缝中心距$d=0.10$ mm、每缝宽$a=0.020$ mm，
波长$\lambda=480$ nm光垂直照射，焦距$f=50$ cm的凸透镜
将远场图样形成于焦面。求：(1)单缝衍射中央明纹宽度；
(2)其包络内可以出现几条双缝干涉主亮纹。
\workings
单缝中央包络由两侧第一暗纹界定，
\[
\boxed{W=\frac{2f\lambda}{a}
=\frac{2(0.50)(480\times10^{-9})}
{0.020\times10^{-3}}
=0.024\,\mathrm m=24\,\mathrm{mm}}.
\]
双缝干涉亮纹间隔
$\Delta x=f\lambda/d=2.40$ mm。
在中央包络内部，
光栅亮纹级数满足
$|m|\lambda/d<\lambda/a$，
即$|m|<d/a=5$。
故$m=-4,-3,\ldots,3,4$，
\[
\boxed{\text{中央单缝包络内共有9条双缝亮纹}}.
\]
\crosscheck $m=\pm5$恰落在单缝第一暗纹位置，
由于\emph{缺级}不计入亮纹，不能算作11条。

\sourceproblem{10.3.19}{274}{274}
\problemtext 光强$I_0$的自然光穿过两片理想偏振片，
初始透振方向平行。保持第一片不动，
将第二片的透振轴分别转过$30^\circ$、$45^\circ$、$60^\circ$，
求每种情况下的最终光强与$I_0$之比。
\workings
第一片使自然光强度成为$I_0/2$。
第二片依马吕斯定律
$I(\theta)=\tfrac12I_0\cos^2\theta$，从而
\[
\boxed{\frac{I_{30}}{I_0}=\frac38},\qquad
\boxed{\frac{I_{45}}{I_0}=\frac14},\qquad
\boxed{\frac{I_{60}}{I_0}=\frac18}.
\]
\crosscheck 透过第一片时必先有$1/2$损耗；
三个出射强度随两轴夹角增大而单调减小。

\sourceproblem{10.3.20}{274}{274}
\problemtext 自然光经过两片透振轴夹角$60^\circ$的偏振片，
输出光强$I_1$；在两片之间插入第三片，其透振轴与前两片
各成$30^\circ$，求新的输出强度$I$与$I_1$之比。
\workings
原两片给
$I_1=(I_0/2)\cos^260^\circ=I_0/8$。
插片后有两段各$30^\circ$的马吕斯衰减，
\[
I=\frac{I_0}{2}\cos^230^\circ\cos^230^\circ
=\frac{I_0}{2}\left(\frac34\right)^2
=\frac{9I_0}{32}.
\]
故
\[
\boxed{\frac I{I_1}=\frac{9/32}{1/8}=\frac94}.
\]
\crosscheck 插片会使出射光强\emph{提高}，
这是逐步转动偏振方向的经典例子，
不违反能量守恒，因为中间片仍在吸收其未透射的分量。

\sourceproblem{10.3.21}{274}{274}
\problemtext 自然光入射两片理想偏振片，求两透振方向夹角$\theta$，
使得通过两片后的光强：
(1)是该两片系统所能达到\emph{最大透射光强}的$1/3$；
(2)是\emph{原入射自然光强度}的$1/3$。
\workings
总透射强度
$I(\theta)=I_0\cos^2\theta/2$，
该系统最大透射光强$I_{\max}=I_0/2$。
(1) $I/I_{\max}=1/3$，
$\cos^2\theta=1/3$，
\[
\boxed{\theta_1=\arccos(1/\sqrt3)\approx54.74^\circ}.
\]
(2) $I/I_0=1/3$，
$\cos^2\theta=2/3$，
\[
\boxed{\theta_2=\arccos\sqrt{2/3}\approx35.26^\circ}.
\]
\crosscheck 题目的两个分母不同：一个是\emph{两片系统的最大透射强度}，
另一个是\emph{自然光入射强度}，所以结果不相同。

\sourceproblem{10.3.22}{274}{274}
\problemtext 一束自然光由空气入射折射率$n=1.40$的液体表面，
若反射光为完全线偏振光，求入射角及折射角。
\workings
反射光完全偏振的布儒斯特条件
$\tan i_B=n_2/n_1=1.40$，故
\[
\boxed{i_B=\arctan1.40\approx54.46^\circ}.
\]
布儒斯特入射时反射线与折射线互相垂直，
\[
\boxed{r_B=90^\circ-i_B\approx35.54^\circ}.
\]
\crosscheck 折射光只是\emph{部分偏振}；
题设“反射完全线偏振”不等于所有折射光也完全偏振。

\sourceproblem{10.3.23}{274}{274}
\problemtext 光从空气射向折射率$n$的玻璃，原书图10.3.23列出
不同的入射光偏振标记与入射角$i$、$i_0=\arctan n$
共五种情形；要求在原图补标反射与折射光的振动方向，
并分别区分线偏振与部分偏振。
\workings
以入射面为参考，电场垂直入射面的分量称$s$偏振，
平行入射面的分量称$p$偏振。
非磁各向同性介质界面的菲涅耳反射/透射不将$s,p$互相转换。
因此可按以下规则给每条射线添加标记：
\[
\boxed{\begin{array}{c|c|c}
\text{入射类型}&\text{反射光}&\text{折射光}\\ \hline
s\text{线偏振}&s\text{线偏振}&s\text{线偏振}\\
p\text{线偏振},\ i\ne i_0&p\text{线偏振}&p\text{线偏振}\\
p\text{线偏振},\ i=i_0&\text{无反射光}&p\text{线偏振}\\
\text{自然光},\ i\ne i_0&\text{部分偏振}&\text{部分偏振}\\
\text{自然光},\ i=i_0&s\text{完全线偏振}&\text{部分偏振}
\end{array}}
\]
布儒斯特角对应$p$分量的反射振幅系数$r_p=0$，
而$s$反射一般不为零。图中“点”与“短划线”必须根据
\emph{纸面是不是入射面}的标识与这些规则匹配。
\crosscheck \textbf{保留原图校核标记}：
源图五个子图的逐光线点线标注不能仅凭文字代替，
此处给出不依赖绘图风格的物理分类表；
后续须在LaTeX中重画并逐图对齐，才可认为原题图示部分完整录入。
